\scanpage{55}
\mspage{4}

Pour $2'$, c'est le cor. 2 ci-dessus, \deleted{\unsure{Pour $3'$ c'est}}\inserted{pour $3'$, cela} résulte aussitôt de la prop. 2. (NB Cet énoncé resterait valable, si on avait pris \textcircled{2} sous la forme faible seulement, cf le NB page précédente)

Je dois encore mettre un axiome de la forme L5, L6 (cf $\overline{\underline{\mathrm{VI}}}$, p. 43-44). Mais l'énoncé de ces axiomes utilise la notion de $\mathrm{Hot}$-\unsure{fibrés}, ou de \deleted{\unsure{morphismes}} foncteurs parfaits, i.e. fait appel à la connaissance de $W_\infty$, ce que j'aimerais éviter (si je peux). Je dois donc introduire la notion de $W$-fibration, sur le modèle de ce que j'ai fait dans le contexte des dérivateurs, dans I (cf p. 111-\unsure{126|125}). On introduit la notion de $\mathbb{D}$-fibration à la p\unsure{age} 117 (\unsure{prop} 5), \unsure{quatre|par} conditions équivalentes, dont certaines ont un sens en termes de la donnée de $W$ sans plus. Soit donc \deleted{$f$}

\[
f:X\unsure{\text{---}|\longrightarrow}Y
\]
\deleted{\unsure{a)}} dans $\mathrm{Cat}$, et considérons les conditions

a) Si $Y''\longrightarrow Y'$ dans $\mathrm{Cat}_{/Y}$ est dans $W$, alors $X''\longrightarrow X'$ aussi (le chgt de base par $X\longrightarrow Y$ transforme $W$ dans $W$)

a$'$) Comme a), mais avec $Y'',Y'$ $W$-asphériques.

$a_{\deleted{\unsure{1}}1}$) Comme a), dans le cas particulier où $Y'$ a un objet final \deleted{\unsure{et}} et $Y''$ \deleted{\unsure{celle}} est la \deleted{\unsure{catégorie de}}\inserted{catégorie ponctuelle} \deleted{\unsure{admise}}

$\deleted{\unsure{a}}a_2$) Dual de $a_1$ avec objet initial de $Y'$.

On a \deleted{\unsure{donc}}
\diagram{\begin{tikzcd}[ampersand replacement=\&,row sep=small,column sep=tiny]
\& a\quad W\text{-fibration }\deleted{\text{\unsure{faible}}}\text{\unsure{ultra}forte} \arrow[d,Rightarrow] \&\\
\& a'\quad W\text{-fibration forte} \arrow[dl,Rightarrow] \arrow[dr,Rightarrow] \&\\
W\text{-fibration gauche}\ a_1 \& \& a_2\quad W\text{-fibration droite}
\end{tikzcd}}

et dans le cas $W=W_\infty$ toutes ces conditions sont équivalentes. J'ai \deleted{\unsure{envie}}\inserted{envie} prendre la plus forte de ces conditions, mais je dois me donner des moyens pour trouver des $W$-fibrations en un sens assez fort. Par exemple, j'aurais \unsure{envie} de savoir que si $X$

% archive sheet number: 55
